\chapter{\ifenglish Conclusions and Discussions\else บทสรุปและข้อเสนอแนะ\fi}

โครงงานวิจัยการประเมินความเป็นส่วนตัวข้อมูลในบริบทภาษาไทยสำหรับตัวแบบภาษาขนาดใหญ่ ได้ดำเนินการทดสอบและวิเคราะห์ผลอย่างครอบคลุม ซึ่งสามารถสรุปผลการดำเนินงาน ปัญหาที่พบ และแนวทางการพัฒนาต่อได้ดังนี้

\section{\ifenglish Conclusions\else สรุปผล\fi}
\begin{sloppypar}
จากการทดสอบโมเดล llama-3.1-8b-instant (Llama), meta-llama-llama-4-maverick-17b-128e-instruct (Meta) และ moonshotai-kimi-k2-instruct-0905 (Kimi) ด้วยการโจมตีรูปแบบ DEA และ JA ในบริบทภาษาไทย 5 ระดับและภาษาอังกฤษ พบว่า:
\begin{itemize}
\item \textbf{ประสิทธิภาพของโมเดล:} llama-3.1-8b-instant มีความโดดเด่นในการต้านทานการโจมตีด้วยการสกัดข้อมูล (Data Extraction Attack หรือ DEA) เนื่องจากมีความสามารถในการกักเก็บข้อมูล (Data Retention) ได้ดีที่สุด ในขณะที่ meta-llama-llama-4-maverick-17b-128e-instruct แสดงความแข็งแกร่งสูงสุดในการยับยั้งการเจลเบรก (Jailbreak Attack หรือ JA) เนื่องจากเป็นโมเดลขนาดใหญ่กว่าที่ใช้หลักการ Semantic Understanding จึงสามารถตรวจจับเจตนาโจมตีได้ดีกว่า

\item \textbf{อิทธิพลของภาษาและระดับภาษาไทย:} ภาษาอังกฤษมีอัตราการโจมตีสำเร็จสูงกว่าภาษาไทยอย่างชัดเจนทั้งในการโจมตีแบบ DEA และ JA เนื่องจากภาษาอังกฤษใช้กระบวนการ Tokenization แบบ subword ซึ่งเอื้อต่อความเข้าใจเชิงความหมายและโครงสร้างประโยคที่ซับซ้อน ในขณะที่ภาษาไทยแบ่งข้อความออกเป็นตัวอักษรเดี่ยวหรือ ASCII รวมถึงการมีข้อมูลสำหรับ train ไม่เพียงพอ สำหรับบริบทภาษาไทย ระดับภาษากันเอง (Casual) มีความทนทานต่อการถูกโจมตีแบบ JA ได้ดีที่สุด เนื่องจาก Training Data ส่วนใหญ่ที่ใช้ในการ Train LLM เป็นข้อมูลที่เป็นภาษาระดับ Casual ซึ่งทำให้ Safety Alignment ป้องกันได้ดีกว่า ในทางกลับกัน ภาษาระดับทางการ (Formal) กลับมีอัตราการโจมตีสำเร็จเฉลี่ยสูงสุด เนื่องจากโครงสร้างภาษาที่เป็นทางการช่วยเสริมความน่าเชื่อถือให้ Prompt โจมตี

\item \textbf{ผลกระทบของความยาวข้อความ (Token Length):} ความยาวของ Prompt มีผลอย่างมากต่อความสำเร็จของการโจมตี DEA โดยโมเดลจะมีจุดอ่อนในช่วงความยาว 5,000--7,000 tokens ซึ่งอัตราการโจมตีสำเร็จจะพุ่งสูงสุด เนื่องจากปรากฏการณ์ Lost-in-the-Middle ที่ทำให้ System Prompt ถูกละเลยเมื่ออยู่กลาง Context ยาว และ Attention Dilution ที่ทำให้ Attention Weights กระจายตัวไปยัง Token จำนวนมาก จน Guard Rails ของโมเดลอ่อนแอลง

\item \textbf{ประสิทธิผลของกลไกการป้องกัน:}
\begin{itemize}
    \item \textbf{Input Scrubbing:} สามารถลดอัตราความสำเร็จของ DEA ลงได้จนใกล้เคียง 0\% แต่ยังมีข้อจำกัดอย่างมากเนื่องจากกรณีชื่อที่มีรูปแบบพิเศษ การแยกชื่อ-นามสกุลด้วยการเว้นบรรทัด และชื่อทับศัพท์ภาษาไทยที่ NER ตรวจจับได้ยาก สำหรับการป้องกัน JA นั้น Scrubbing จะได้ผลดีกับโมเดลขนาดเล็ก (llama) เนื่องจากพึ่งพา Surface-level pattern matching แต่แทบไม่ได้ผลกับโมเดล Meta และ Kimi เนื่องจากเป็นโมเดลที่ใช้หลักการ Semantic Understanding
    \item \textbf{Defensive Prompting:} สามารถลดอัตราความสำเร็จของ DEA ได้เข้าใกล้ 0\% เช่นกัน เนื่องจากกลไกการสวมบทบาท (Role-Playing) และการดักจับด้วยเงื่อนไขที่ตรงกัน นอกจากนี้ยังสามารถลดอัตราการโจมตีแบบ JA ได้ระดับหนึ่ง แต่ยังคงมีความเปราะบางอยู่บ้าง เนื่องจาก Jailbreak Prompt มีความหลากหลายทางโครงสร้างมากกว่า DEA ทำให้ Defensive Prompt ไม่สามารถดักจับได้ทุกรูปแบบ
\end{itemize}
\end{itemize}
\end{sloppypar}

ข้อจำกัดของระบบ:
\begin{itemize}
    \item ความแม่นยำของภาษาไทย: ระบบ Scrubbing ยังมีข้อจำกัดในการตรวจจับชื่อบุคคลไทยที่ไม่มีรูปแบบตายตัว หรือคำแสลงในระดับภาษากันเอง (Casual)
    \item ข้อจำกัดด้านทรัพยากร: การทดสอบถูกจำกัดด้วยอัตราการเรียกใช้งาน (Rate Limit) ของ Groq API และงบประมาณในการเข้าถึงทรัพยากรทางเทคนิค
    \item ขอบเขตข้อมูล: ผลการทดลองอ้างอิงจากชุดข้อมูลอีเมล Enron และข้อมูลบุคคลสำคัญ ซึ่งอาจไม่ครอบคลุมข้อมูลส่วนบุคคลในบริบทอื่น ๆ ทั้งหมด
\end{itemize}

\section{\ifenglish Challenges\else ปัญหาที่พบและแนวทางการแก้ไข\fi}

ในการทำโครงงานนี้ คณะผู้จัดทำพบปัญหาหลักและได้ดำเนินการแก้ไขดังนี้:
\begin{enumerate}
    \item ปัญหาด้านเสถียรภาพของ API และข้อจำกัดความเร็ว
        \begin{itemize}
            \item ปัญหา: การส่งคำสั่งโจมตีปริมาณมากทำให้เกิด Error และการหยุดชะงักของระบบ
            \item การแก้ไข: พัฒนาระบบรันอัตโนมัติที่รองรับหลาย API Key และเพิ่มระบบ Checkpoint เพื่อให้สามารถรันงานต่อจากจุดที่ค้างได้ทันที
        \end{itemize}
    \item ความซับซ้อนในการวัดผลการโจมตี (Metrics Accuracy)
        \begin{itemize}
            \item ปัญหา: มาตรวัด JA ในระยะแรกไม่สามารถแยกแยะการตอบกลับที่เป็นอันตรายได้อย่างแม่นยำ
            \item การแก้ไข: ปรับปรุงมาตรวัด (Edit Metrics) โดยเน้นการตรวจสอบ Contact Information และวิเคราะห์ Keyword ร่วมกันอย่างละเอียด
        \end{itemize}
    \item ความคลาดเคลื่อนในการแปลระดับภาษา
        \begin{itemize}
            \item การแปล Prompt ปริมาณมากด้วยระบบอัตโนมัติอาจทำให้โทนของระดับภาษาผิดเพี้ยน
            \item การแก้ไข: ใช้การประมวลผลแบบ Batch ผ่าน OpenAI API และทำการตรวจสอบความถูกต้อง (Manual Verification) ในบางส่วนเพื่อให้มั่นใจว่าครอบคลุมทั้ง 5 ระดับภาษาจริง
        \end{itemize}
    \item การจัดการผลลัพธ์จำนวนมาก
        \begin{itemize}
            \item ปัญหา: ข้อมูลผลลัพธ์มีปริมาณมากจนยากต่อการวิเคราะห์ด้วยมือ
            \item การแก้ไข: สร้างสคริปต์อัตโนมัติเพื่อคำนวณ Token, ความยาวข้อความ และสร้างกราฟสรุปผลลงใน Google Sheets โดยตรง
        \end{itemize}
\end{enumerate}
\newpage
\section{\ifenglish%
Suggestions and further improvements
\else%
ข้อเสนอแนะและแนวทางการพัฒนาต่อ
\fi
}

เพื่อให้งานวิจัยนี้มีความสมบูรณ์และเป็นประโยชน์ต่อวงการ AI ในประเทศไทยยิ่งขึ้น มีข้อเสนอแนะดังนี้:
\begin{itemize}
    \item การขยายขอบเขตโมเดล: ควรมีการทดสอบกับโมเดลที่ถูกปรับจูนมาเพื่อภาษาไทยโดยเฉพาะ (Thai Fine-tuned Models) เพื่อเปรียบเทียบความปลอดภัยกับโมเดลมาตรฐานสากล
    \item การพัฒนาระบบ Scrubbing เชิงลึก: พัฒนาโมเดล NER ที่มีความเฉพาะเจาะจงกับบริบทของประเทศไทยมากขึ้น เพื่อเพิ่มความแม่นยำในการคัดกรองข้อมูลก่อนส่งเข้าสู่ LLMs
    \item การทดสอบในสภาพแวดล้อมจริง: นำกลไกการป้องกันที่ได้จากงานวิจัยไปทดสอบกับแอปพลิเคชันที่ใช้งานจริง เช่น ระบบ Customer Service เพื่อประเมินผลกระทบต่อความล่าช้า (Latency) และประสบการณ์ของผู้ใช้งาน
    \item การปรับปรุง Defensive Prompting อัตโนมัติ: ศึกษาการใช้ AI ในการสร้างและปรับเปลี่ยนคำสั่งป้องกัน (Self-adaptive Defense) ตามรูปแบบการโจมตีที่เปลี่ยนไปในอนาคต
\end{itemize}




